\begin{proofbox}

	\textbf{\textcolor{CProof}{思路提示}}

	由正方体的中心对称性，内切球、棱切球、外接球的球心都与正方体的中心重合，因此只需要计算中心到相应切触元素的距离。

	\medskip
	\textbf{\textcolor{CProof}{正式推导}}
	\begin{description}[style=nextline, leftmargin=2.8em, labelsep=0.5em, itemsep=0.55em]
		\item[\textbf{步骤一（确定球心）：}] 正方体的中心是几何中心，记为$O$。由对称性，若球与正方体各面相切或与各棱相切或过各顶点，则球心必须位于$O$。
			\par\textbf{理由：}
			正方体是中心对称图形，且切接关系具有对称性，只有中心满足到所有面、所有棱中点、所有顶点等距的要求。

		\item[\textbf{步骤二（内切球半径）：}] 内切球与六个面都相切，球心$O$到每个面的距离即为半径。设棱长为$a$，则$O$到任意面的距离为$\frac{a}{2}$。
			\[
				r_i = \frac{a}{2}.
			\]
			\par\textbf{理由：}
			面对面的距离为$a$，中心到任一面距离为$a/2$。

		\item[\textbf{步骤三（棱切球半径）：}] 棱切球与十二条棱的中点相切。取一条棱中点$M$，则$OM$的距离即为半径。由勾股定理，
			\[
				\begin{aligned}
					OM &= \sqrt{\left(\frac{a}{2}\right)^2 + \left(\frac{a}{2}\right)^2} \\
					&= \frac{\sqrt{2}}{2}a.
			\end{aligned}
			\]
			故
			\[
				r_m = \frac{\sqrt{2}}{2}a.
			\]
			\par\textbf{理由：}
			棱中点在面对角线方向偏离中心两个方向各$a/2$，距离由勾股定理得。

		\item[\textbf{步骤四（外接球半径）：}] 外接球过八个顶点。取任一顶点$P$，$OP$的距离即为半径。$P$在三个方向上都偏离中心$a/2$，所以
			\[
				\begin{aligned}
					OP &= \sqrt{\left(\frac{a}{2}\right)^2 + \left(\frac{a}{2}\right)^2 + \left(\frac{a}{2}\right)^2} \\
					&= \frac{\sqrt{3}}{2}a.
			\end{aligned}
			\]
			因此
			\[
				r_o = \frac{\sqrt{3}}{2}a.
			\]
			\par\textbf{理由：}
			顶点坐标三个分量绝对值均为$a/2$，距离公式计算即得。
	\end{description}

	\medskip
	\textbf{\textcolor{CProof}{结论回扣}}

	综上，正方体的内切球、棱切球、外接球半径分别为$\frac{a}{2}$、$\frac{\sqrt{2}}{2}a$、$\frac{\sqrt{3}}{2}a$。这三个公式在立体几何切接问题中可以直接使用。

\end{proofbox}
